\documentstyle[%


righttag%


,11pt%


,amssymb%


,russian%               remove in Eng.


,izvru%                 Set izven% in Eng.


]{amsart}





%


\newcommand{\dist}{\operatorname{dist}}


\newcommand{\supp}{\operatorname{supp}}


\newcommand{\la}{\langle}


\newcommand{\ra}{\rangle}


\newcommand{\tts}[1]{}





\theoremstyle{plain}


\theorembodyfont{\it}


\newtheorem{thms}{\hskip\parindent Теорема}[section]


\newtheorem{lems}{\hskip\parindent Лемма}[section]





\theoremstyle{definition}


\theorembodyfont{\rm}


\newtheorem{cornonum}{\hskip\parindent Следствие}


\newtheorem{notenonum}{\hskip\parindent Замечание}


\newtheorem{defnnonum}{\hskip\parindent Определение}


\renewcommand{\thedefnnonum}{}


\renewcommand{\thenotenonum}{}


\renewcommand{\thecornonum}{}





%\hoffset=-15mm%                  For Eng.


\setcounter{page}{46}


\god{2003}


\nomer{9~(496)} %                        Remove in Eng.


%\nomer{Vol.\,47, No.\,9}%               Set in Eng.


%\str{\pageref{begin}--\pageref{end}}%  Set in Eng.


\udk{517.98}





\author{\it М.Р.\,Тимербаев}


% NB:   ^^ remove in Eng.


\title{Пространства с нормой графика \\ и усиленные пространства


Соболева. II}


\thanks{Работа выполнена при финансовой


поддержке Российского фонда фундаментальных исследований (проекты 


01-01-00616, 03-01-00380) и гранта \No\,Е02-1.0-189 


Министерства образования России по фундаментальным исследованиям


в области естественных и точных наук.}





\date{}


%\pagestyle{myheadings}%         Set in Eng.


\pagestyle{plain} %              Remove in Eng.


\begin{document}


%\thispagestyle{plain}%          Set in Eng.


\maketitle


\label{begin}











\section{Пространства вектор-функций со значениями\protect\\ в


пространствах с нормой графика}








Данная работа является продолжением статьи \cite{Timer}.


Напомним основные понятия и обозначения. Через $\Phi$, $\Psi$


обозначаются топологические векторные пространства (ТВП). Всюду,


если не оговорено особо, для ТВП $U$ включение $U \subset \Phi$


будет пониматься не только в теоретико-множественном, но и в


топологическом смысле; таким образом, равенство $U=V$ для


$B$-пространств будет означать также эквивалентность норм этих


пространств. Отметим также следующее: если $B$-пространства $U, V


\subset \Phi$ таковы, что $V$ является подмножеством $U$, то $V$


непрерывно вложено в $U$ --- это следует из теоремы Банаха о


замкнутом графике, которую нужно применить к тождественному


отображению.





Обозначим через $L(\Phi, \Psi)$ множество


линейных непрерывных отображений из $\Phi$ в $\Psi$. Пусть $\gamma


\in L(\Phi, \Psi)$ и $U \subset \Phi$. Тогда, очевидно, $\gamma


\in L(U, \Psi)$ (сужение $\gamma$ на $U$ здесь и далее будет


обозначаться тем же символом $\gamma$). Если $U \subset \Phi$ ---


$B$-пространство, то линейное множество $\gamma(U)$ (образ $U$ при


отображении $\gamma$), наделенное фактор-нормой


$$


{\|x\|}_{\gamma(U)} = \inf \{ {\|u\|}_U : u \in U,\ \gamma u=x \},


$$


будет $B$-пространством, изометричным, очевидно,


фактор-пространству $U/(\ker \gamma \cap U)$, причем $\gamma(U)


\subset \Psi$ и $\gamma \in L(U, \gamma(U))$. Кроме того, если


$B$-пространство $X$ непрерывно вложено в $\Psi$, то, для того


чтобы $\gamma$ непрерывно отображало $U$ в $X$, достаточно (и,


разумеется, необходимо), чтобы $\gamma(U)$ было подмножеством


пространства $X$. Если $\gamma$ непрерывно отображает


$B$-пространство $U$ на $B$-пространство $X$, то фактор-норма


пространства $\gamma (U)$ эквивалентна норме пространства $X$.





Оператор $\gamma \in L(U, X)$ называется ретракцией, если


существует оператор $\beta \in L(X, U)$, называемый коретракцией,


такой, что


$$


\gamma \beta x = x \quad \forall x \in X.


$$





Пусть $\gamma \in L(\Phi, \Psi)$


и пусть $U \subset \Phi$ и $X \subset


\Psi$ --- два $B$-пространства. Определим пространство


$$


\gamma (U, X) \overset{\mathrm{def}}{=}{} \{u \in U: \gamma u \in X\},


$$


наделенное нормой графика


\begin{equation}


{\|u\|}_{\gamma(U,X)} = {\|u\|}_U + {\|\gamma u\|}_X.


\label{normaV}


\end{equation}


Из определения нормы (\ref{normaV}) и из полноты


пространств $U$, $X$ с очевидностью вытекает, что $\gamma(U,X)$


является $B$-пространством, непрерывно вложенным в пространство


$U$, причем $\gamma (U,X) = U$ тогда и только тогда, когда


$\gamma(U)$ является подмножеством пространства $X$.





Пусть $(T, \Sigma, \mu)$ --- пространство с (полной конечной) мерой


$\mu$, заданной на $\sigma$-алгебре $\Sigma$ подмножеств


множества $T$. Для $B$-пространства $U \subset \Phi$ множество


всех $\mu$-измеримых функций $u:T \to U$ обозначается через


$L_0(T, \Sigma, \mu; U)=L(T;U)$, или, короче, когда не возникает


двусмысленности, через $L_0(U)$. Для простоты изложения


ограничимся рассмотрением сепарабельных $B$-пространств $U


\subset \Phi, X \subset \Psi$. Поэтому в силу критерия Петтиса


измеримости функции (\cite{Danford}, c.\,166) функция $u:T \to U$


будет измерима тогда и только тогда, когда она $U^*$-скалярно


измерима, т.\,е. для каждого элемента $u^* \in U^*$ измерима


скалярная функция $t \to {\langle u(t),u^*\rangle}_U$.





Пусть $M$ --- замкнутое подпространство пространства $U$.


Функцию $t \in T \to f(t) \in U/M$ со значениями в


фактор-пространстве можно рассматривать и как многозначную


функцию $t\in T \to f(t) \in 2^U$. Как показано ниже,


существование измеримого сечения $u(t) \in f(t)$


равносильно измеримости функции $f$.





\begin{lems} \label{lmeas}


Для функции $f:T \to U/M$ эквивалентны следующие утверждения\/${:}$


\begin{itemize}


\item[(i)] функция $f:T \to U/M$ измерима$;$


\item[(ii)] для любого $u \in U$ скалярная функция $t \in T \to


\dist(u, f(t)) = \inf \{ {\|u-v\|}_U: v \in f(t) \}$ измерима$;$


\item[(iii)] для любого $u \in U$ множество $\{t \in T: u \in


f(t)\}$ \ $\mu$-измеримо$;$


\item[(iv)] существует счетное множество измеримых функций


$u_n:T \to U$ таких, что $u_n(t) \in f(t)$ почти всюду


и при почти всех $($п.\,в.$)$ $t


\in T$ множество $\{u_n(t): n \ge 1\}$ плотно в $f(t);$


\item[(v)] существует измеримая функция $u:T \to U$ такая, что


$u(t) \in f(t)$.


\end{itemize}


\end{lems}





\begin{pf}


(i) $\Rightarrow$ (ii). Пусть $u \in U$ и


$\wh {u} = u+M \in U/M$. Если $f:T \to U/M$


измерима, то, очевидно, измерима и норма разности $t \in T \to


{\|\wh {u}-f(t)\|}_{U/M}$. Осталось заметить, что


$$


{\|\wh {u}-f(t)\|}_{U/M}=\dist(u,f(t)).


$$





Эквивалентность утверждений (ii), (iii) и (iv) содержится в


теореме Дебре--Кастена об измеримости многозначных отображений


(\cite{IoffeLevi}, c.\,74). Импликация (iv)~$\Rightarrow$~(v)


тривиальна. Осталось показать (v)~$\Rightarrow$~(i).





Пусть существует измеримое сечение $u(t) \in f(t)$.


Поскольку $B$-пространство $U$ сепарабельно, то сепарабельно и


фактор-пространство $U/M$. По теореме Петтиса об измеримости


достаточно установить, что для любого линейного непрерывного


функционала $F$ на $U/M$ измерима скалярная функция $t\in T \to


{\la f(t),F\ra}_{U/M}$. Но каждый такой функционал $F \in (U/M)^*$


формулой $u \in U \to F(\wh {u})$ определяет


линейный непрерывный функционал $u^* \in U^*$ такой, что $M


\subset \ker u^*$. Поэтому


$$


{\la f(t),F\ra }_{U/M}={\la \wh {u}(t),F\ra}_{U/M}=


{\la u(t),u^*\ra}_U,


$$


последняя же функция измерима в силу измеримости функции $t \to


u(t)$.


\end{pf}





\begin{lems} \label{ueps}


Пусть $f:T \to U/M$ измерима.


Тогда для любого $\varepsilon > 0$ найдется измеримая функция


$u_{\varepsilon}:T \to U$ такая, что $u_{\varepsilon}(t) \in f(t)$ и


$$


{\|f(t)\|}_{U/M} \le {\|u_{\varepsilon}(t)\|}_U \le


(1+\varepsilon){\|f(t)\|}_{U/M} \quad \text{п.\,в.}


$$


\end{lems}





\begin{pf}


Согласно лемме \ref{lmeas} (утверждение (iv))


существует счетное множество измеримых функций $u_n:T \to U$


таких, что $u_n(t) \in f(t)$ почти всюду и при п.\,в. $t \in T$ множество


$\{u_n(t): n \ge 1\}$ плотно в $f(t)$.


Тогда при п.в. $t \in T$ \ ${\|f(t)\|}_{U/M}=


\inf\{ {\|u_n(t)\|}_U : n \ge 1 \}$.


Положим


$T_0 = \{t \in T: 0 \in f(t)\}$


и при фиксированном


$\varepsilon > 0$ \ $T_n = \{t \in T: {\|u_n(t)\|}_U \le


(1+\varepsilon){\|f(t)\|}_{U/M} \}$.


Построенные множества $T_n \ (n \ge 0)$ измеримы и


$\mu \Big(T \setminus \bigcup \limits_{n=0}^{\infty} T_n\Big)=0$. Далее


строим последовательность непересекающихся множеств $E_n$,


полагая при $n \ge 1$ \ $E_n = T_n \setminus \bigcup \limits_{j=0}^{n-1}


T_j$. Функцию $u_{\varepsilon}$ определим формулой


$$


u_{\varepsilon}(t) = \sum \limits_{n=1}^{\infty}


{\chi}_{_{E_n}}(t) u_n(t),


$$


где ${\chi}_{_{E_n}}(t)$ обозначает характеристическую функцию


множества $E_n$. По построению $u_{\varepsilon}(t) \in f(t)$ при


п.\,в. $t \in T$, а потому ${\|f(t)\|}_{U/M} \le


{\|u_{\varepsilon}(t)\|}_U$. Кроме того, как следует из


определения множеств $T_0$ и $E_n$,


${\|u_{\varepsilon}(t)\|}_U \le (1+\varepsilon){\|f(t)\|}_{U/M}$.


\end{pf}








Пусть $L_p = L_p(T, \Sigma, \mu)$. Как обычно, пространство


$U$-значных функций класса $L_p$ ($1 \le p < \infty$)


c соответствующей метрикой обозначаются


через $L_p(T, \Sigma, \mu; U)=L_p(T;U)$ (коротко


$L_p(U)$). Если $\gamma \in L(\Phi, \Psi)$ и


$U \subset \Phi$, $X \subset \Psi $,


то для любой функции $u \in L_p(U)$ функция


$t \in T \rightarrow \gamma u(t) \in X$ будет


принадлежать классу $L_p(X)$. Тем самым


$\gamma$ определяет отображение из $L_p(U)$


в $L_p(X)$, которое будем обозначать


тем же символом $\gamma$.





\begin{thms} \label{Lp}


Имеет место теоретико-множественное и


метрическое равенство


$$


\gamma (L_p(U)) = L_p(\gamma (U)).


$$


\end{thms}





\begin{pf}


Положим $X=\gamma(U)$. Тогда по определению


${\left\| {\gamma} \right\|}_{U \rightarrow X} =1$, откуда следует, что


${\left\| {\gamma} \right\|}_{L_p(U) \rightarrow L_p(X)} \leq 1$


и


$\gamma (L_p(U)) \subset L_p(X)$. Пусть $x \in L_p(X)$ ---


произвольная функция. Рассматривая $\gamma^{-1}$


как изометрию из $X$ в $U/U_0$, определим функцию


$f(t) = \gamma^{-1}x(t) \in U/U_0$.


По предыдущей лемме для произвольного $\varepsilon > 0$ найдется 


такая измеримая функция $u(t) \in f(t)$, что


${\|f(t)\|}_{U/U_0} \le {\|u(t)\|}_U \le


(1+\varepsilon){\|f(t)\|}_{U/U_0}$ почти всюду.


Итак, для произвольного $\varepsilon$ найдена функция $u$ такая, 


что $x(t) = \gamma u(t)$ почти всюду и


$$


{\left\| {x} \right\|}_{L_p(X)} = {\|f\|}_{L_p(U/U_0)} \le {\|u\|}_{L_p(U)}


\le (1+\varepsilon){\|f\|}_{L_p(U/U_0)} = (1+\varepsilon){\|x\|}_{L_p(X)}.


$$


В силу произвольности $x \in L_p(X)$


и $\varepsilon > 0$ это означает, что множества


$\gamma (L_p(U))$ и $L_p(X)$ совпадают между собой и


${\left\| {x} \right\|}_{L_p(X)} =


{\left\| {x} \right\|}_{\gamma (L_p(U))}$.


\end{pf}





\begin{thms} \label{LpUX}


Имеет место теоретико-множественное и топологическое равенство


$$


\gamma (L_p(U), L_p(X)) = L_p(\gamma (U,X)).


$$


\end{thms}





{\bf Доказательство.}


Введем обозначения: $V = \gamma(U,X)$,


$W = \gamma(L_p(U), L_p(X))$. Если $u \in W$, то $u(t) \in V$ почти всюду.


Пусть $c_1, c_2 > 0$ такие постоянные (зависящие от $p$),


для которых выполняются неравенства


$$


c_1(|a|^p + |b|^p) \le (|a|+|b|)^p \le c_2(|a|^p + |b|^p)


$$


($a$, $b$ --- произвольные числа). Тогда для произвольной 


функции $u \in W \cup L_p(V)$


справедлива цепочка неравенств


$c_1{\left\| {u} \right\|}_W^{p} \le


{\left\| {u} \right\|}_{L_p(V)}^p \le c_2 {\left\| {u} \right\|}_W^{p}$,


откуда вытекает утверждение.


\medskip








Для $1 \le p,q < \infty$ рассмотрим теперь усиленное пространство


$W=\gamma(L_p(U), L_q(X))$. Ясно, что если $u \in W$,


то $u(t) \in \gamma(U,X)$


при п.в. $t \in T$.





\begin{lems} \label{ldensW}


Множество простых $($т.\,е. измеримых и принимающих конечное число


зна\-че\-ний$)$ $\gamma(U,X)$-знач\-ных функций плотно в $W$.


\end{lems}





{\bf Доказательство} этого утверждения проводится по той же схеме, что и для


пространств типа $L_p$. Cначала докажем плотность в усиленном


пространстве $W$ подмножества счетнозначных измеримых функций из $W$:


действительно, для любой функции $u \in W$ и произвольного $\varepsilon >0$


найдется счетнозначная измеримая функция $u_{\varepsilon}$ такая, что при


п.\,в. $t \in T$ \ ${\left\| {u(t)-u_{\varepsilon}(t)}


\right\|}_{\gamma(U,X)} < \varepsilon$,


отсюда $u_{\varepsilon} \in W$ и


${\left\| {u - u_{\varepsilon}} \right\|}_W < c \varepsilon$


($c$ --- постоянная, не зависящая от $u$ и $\varepsilon$).


Далее, т.\,к. норма пространства


$W$ абсолютно непрерывна, т.\,е.


$$


{\left\| {\chi_{_E} u} \right\| }_W \rightarrow 0


\quad \text{при} \quad \mu(E) \rightarrow 0,


$$


найдется такое множество $E \in \Sigma$, что ``срезка''


${\chi}_{_E} u_\varepsilon $ функции $u_\varepsilon$ будет конечнозначной и


${\left\| {u - {\chi}_{_E}u_{\varepsilon}} \right\|}_W < c \varepsilon$,


что доказывает утверждение.





\begin{thms} \label{dens}


Пусть множество скалярных функций $L \subset L_p \cap L_q$ плотно


в $L_p \cap L_q$ и множество $M \subset \gamma(U,X)$ плотно в $\gamma(U,X)$.


Тогда линейная оболочка множества функций


$L \otimes M = \left\{ {\varphi u: \varphi \in L,


\ u \in M} \right\}$ плотна в пространстве $W=


\gamma(L_p(U), L_q(X))$.


\end{thms}





\begin{pf}


Пусть $r=\max (p,q)$, $V=\gamma(U,X)$. Из условия


теоремы вытекает плотность линейной оболочки множества функций


$L \otimes M$ в $L_r(V)$. В свою очередь, из леммы \ref{ldensW}


следует плотность $L_r(V)$ в $W$.


\end{pf}





\begin{cornonum}


Рассмотрим единичный интервал


$T=(0,1)$ с лебеговой мерой. Если множество $M \subset


\gamma(U,X)$ плотно в пространстве $\gamma(U,X)$, то линейная


оболочка множества функций $\{\varphi u:$ $ \varphi \in


C^{\infty}[0,1]$, $u \in M\}$ плотна в пространстве


$W=\gamma(L_p(T;U),$ $ L_q(T;X))$. 


\end{cornonum}








\section{Пространства Соболева с усиленной метрикой}





Пусть $\Omega$ --- ограниченная плоская область 


с границей $\Gamma_0=\partial \Omega$, состоящей из конечного числа 


кусков класса $C^2$, причем соседние куски в точках стыка


не образуют нулевых углов. Пусть даны кривые $\Gamma_i \subset


\overline{\Omega}$, $i=1,2,\dotsc,m$,


класса $C^2$ с концами на $\Gamma_0$. Предполагается, что в точках


пересечений кривые $\Gamma_i$ ($i=\overline{0,m}$)


также не образуют нулевых углов.


Таким образом, область $\Omega$ разбивается на конечное число подобластей


$\Omega_k$, $k=\overline{1,n}$, с границами того же класса, что и у $\Omega$.





\begin{defnnonum}


Пусть $1<p, q<\infty$.


Пространством Соболева с усиленной метрикой называют множество


функций $$


V=V_{p,q}^{1,1}=V_{p,q}^{1,1}(\Omega, \{\Gamma_i\})


\overset{\mathrm{def}}{=}{} \{u \in W_p^1(\Omega):


{u}|_{\Gamma_i} \in W_q^1(\Gamma_i), \ i=\overline{0,m} \}


$$


с нормой


\begin{equation}


{\left\| {u} \right\|}_V=


{\left\| {u} \right\|}_{W_p^1(\Omega)}+\sum \limits_{i=0}^{m}


{\left\| {u} \right\|}_{W_q^1(\Gamma_i)}. \label{norma}


\end{equation}


\end{defnnonum}





Пространства такого рода рассматривались в 


\cite{Anton}--\cite{DjakIzVuz}. Из определения с очевидностью


вытекает, что пространство $V$ является рефлексивным


$B$-пространством. Ниже установим плотность гладких функций в


усиленном пространстве Соболева, а также в одном специальном


пространстве, которое можно интерпретировать как пространство


$V$-значных функций того же типа, что и в предыдущем пункте. Этот


факт позволяет дать эквивалентное определение пространства


Соболева с усиленной метрикой как замыкание гладких функций в


норме (\ref{norma}).





Пусть $\Gamma=\bigcup \limits_{i=0}^{m} \Gamma_i$. Для $r \in (0,1]$ через


$W_p^r(\Gamma)$ обозначим множество функций, определенных на $\Gamma$


таких, что сужение этих функций на каждое $\Gamma_i$ принадлежит классу


$W_p^r(\Gamma_i)$


и, кроме того, в случае $rp>1$ функции класса $W_p^r(\Gamma)$


предполагаются (после исправления на множестве нулевой линейной меры)


непрерывными на компакте $\Gamma$. На $W_p^r(\Gamma)$ естественно ввести


нормировку


\begin{alignat*}{2}


{\left\| {u} \right\|}_{W_p^r(\Gamma)}&=


\sum \limits_{i=0}^{m} {\left\| {u} \right\|}_{W_p^r(\Gamma_i)}


&\quad &\text{при $rp \le 1$} \\


%


\intertext{и}


%


{\left\| {u} \right\|}_{W_p^r(\Gamma)}&=


\sum \limits_{i=0}^{m} {\left\| {u} \right\|}_{W_p^r(\Gamma_i)}


+ {\left\| {u} \right\|}_{C(\Gamma)}


&\quad  &\text{при $rp>1$}.


\end{alignat*}


В последнем случае можно наделить $W_p^r(\Gamma)$ эквивалентной нормой


$$


{\left\| {u} \right\|}_{W_p^r(\Gamma)}=


\sum \limits_{i=0}^{m} {\left\| {u} \right\|}_{W_p^r(\Gamma_i)}


+ \bigg( {\sum \limits_{\xi } |u(\xi)|^p} \bigg)^{1/p},


$$


где точка $\xi$ пробегает все пересечения семейства 


кривых $\{\Gamma_i\}$, $i=\overline{0,m}$.


Очевидно, $W_p^r(\Gamma)$ является рефлексивным


$B$-пространством.





Всюду в дальнейшем через $\gamma$ будет обозначаться оператор


следа на $\Gamma$. Оператор $\gamma$ является ретракцией из


$W_p^1(\Omega)$ на $W_p^{1-1/p}(\Gamma)$. Действительно, поскольку


граница каждой из подобластей $\Omega_k$ состоит из гладких


кусков, стыкующихся без нулевых углов, то, как следует из


результатов работы \cite{Jak}, операторы следа $\gamma_k u = {u


}|_{\partial \Omega_k}$ являются ретракциями из $W_p^1(\Omega_k)$


на $W_p^{1-1/p}(\partial \Omega_k)$. Поэтому существуют


соответствующие коретракции (операторы продолжения) $\beta_k \in


L(W_p^{1-1/p}(\partial \Omega_k), W_p^1(\Omega_k))$. Очевидно,


определяемый ``кусочно'' оператор


$$


\beta \varphi = \sum_{k=1}^{n} \chi_{\Omega_k} \beta_k \delta_k \varphi,


$$


где $\delta_k$ --- операторы сужения с $\Gamma$ 


на $\partial \Omega_k$, будет коретракцией для $\gamma$.








\begin{thms} \label{retr}


Пространство $V$ является


пространством $\gamma (W_p^1(\Omega), W_q^1(\Gamma))$ с нормой


графика в смысле определения п.\,$2$ $\cite{Timer}$, причем оператор


следа $\gamma$ является ретракцией $V$ на $W_q^1(\Gamma)\cap


W_p^{1-1/p}(\Gamma)$.


\end{thms}





\begin{pf}


Включение $\gamma (W_p^1(\Omega),


W_q^1(\Gamma)) \subset V$ очевидно. Для доказательства


противоположного включения покажем, что для любой функции $u \in


V$ ее след $\gamma u$ является непрерывной функцией.


Действительно, поскольку $W_q^1(\Gamma_i) \subset C(\Gamma_i)$, то


${u}|_{\Gamma_i} \in C(\Gamma_i)$ ($i=\overline{0,m}$). Отсюда и


из того, что $\gamma u \in W_p^{1-1/p}(\Gamma)$ следует, что если


$\xi$ --- точка пересечения кривых $\Gamma_i$ и $\Gamma_j$, то в


силу условий ``склейки'' (\cite{Jak}, c.\,74) $\gamma u$ должна быть


непрерывной в точке $\xi$, поэтому $\gamma u \in C(\Gamma)$.





То, что оператор $\gamma$ является ретракцией из $V$ на


$W_q^1(\Gamma)\bigcap W_p^{1-1/p}(\Gamma)$, следует из того, что


$\gamma$ есть ретракция из $W_p^1(\Omega)$ на


$W_p^{1-1/p}(\Gamma)$ и из утверждения (iii) теоремы 2.1


\cite{Timer}.


\end{pf}





\begin{notenonum}


Как следует из теорем вложения, при выполнении неравенства


$2q \geq p$ будет выполнено включение $W_q^1(\Gamma) \subset


W_p^{1-1/p}(\Gamma)$. В этом случае $\gamma$ будет являться ретракцией


из $V$ на $W_q^1(\Gamma)$.


\end{notenonum}








Через $\overset{\circ}{V}$ обозначим подпространство функций $u$


из $V$, для которых $\gamma u = 0$. Ясно, что $\overset{\circ}{V}$


является замкнутым подпространством пространств $W_p^1(\Omega)$ и


$V$. Как было отмечено выше, оператор $\gamma$ является ретракцией


из $W_p^1(\Omega)$ на $W_p^{1-1/p}(\Gamma)$ и, как утверждается в


теореме \ref{retr}, из $V$ на $W_q^1(\Gamma)\cap


W_p^{1-1/p}(\Gamma)$. Отсюда вытекает дополняемость


$\overset{\circ}{V}$ как в $W_p^1(\Omega)$, так и в $V$. Пусть


$\pi$ --- какой-нибудь проектор на $\overset{\circ}{V}$. Тогда из


утверждения (iv) теоремы 2.1 \cite{Timer} получаем





\begin{cornonum}


Оператор $ \Pi u = (\pi u, \gamma


u)$ осуществляет изоморфизм пространства $V$ на декартово


произведение $\overset{\circ}{V}{} \times (W_q^1(\Gamma)\cap


W_p^{1-1/p}(\Gamma))$. 


\end{cornonum}








Установим теперь плотность гладких функций в пространствах


Соболева с усиленной метрикой. Пусть $D_{\Gamma}(\Omega)$


обозначает множество функций, имеющих частные производные любого


порядка и носители которых не пересекаются с $\Gamma$. Справедлива





\begin{thms} \label{C0densV0}


Множество $D_{\Gamma}(\Omega)$ плотно в $\overset{\circ}{V}$.


\end{thms}





{\bf Доказательство.}


Пусть $u \in \overset{\circ}{V}$. Тогда $u|_{\Omega_k} \in


\overset{\circ}{W}{}_p^1(\Omega_k)$ для каждого $k=\overline{1,n}$. Поскольку


пространство финитных функций $D(\Omega_k)$ плотно в


$\overset{\circ}{W}{}_p^1(\Omega_k)$


(\cite{Necas}, c.\,87), то для фиксированного


$\varepsilon > 0$ можно подобрать функцию


$u_{\varepsilon, k} \in D(\Omega_k)$ такую, что


${\left\| {u-u_{\varepsilon,k}} \right\|}_{W_p^1(\Omega_k)} \le


\varepsilon/n$.


Положив


$u_\varepsilon=\sum \limits_{k=1}^{n} \chi_{\Omega_k}u_{\varepsilon,k}$,


очевидно, получим $u_\varepsilon \in D_{\Gamma}(\Omega)$ и


$$


{\left\| {u-u_\varepsilon} \right\|}_V =


{\left\| {u-u_\varepsilon} \right\|}_{W_p^1(\Omega)} \le


\sum \limits_{k=1}^{n}{\left\|


{u-u_{\varepsilon,k}} \right\|}_{W_p^1(\Omega_k)} \le \varepsilon.


\quad\square


$$





\begin{thms} \label{CdensV}


Множество бесконечно дифференцируемых функций


$C^\infty(\overline{\Omega})$ плотно в пространстве $V$.


\end{thms}





{\bf Доказательство.}


Положим $r=\max (p,q)$. Хорошо известно, что


множество $C^\infty(\overline{\Omega})$ плотно в $W_r^2(\Omega)$.


Поэтому достаточно установить плотность последнего в усиленном


пространстве Соболева $V$ (заметим, что $W_r^2(\Omega) \subset


V$). Согласно следствию 3 теоремы 2.1 \cite{Timer} для этого достаточно


установить, что множество $W_r^2(\Omega) \bigcap \ker \gamma$


плотно в $\overset{\circ}{V}$ и $\gamma(W_r^2(\Omega))$ плотно в


$\gamma(V)=W_q^1(\Gamma) \bigcap W_p^{1-1/p}(\Gamma)$. Поскольку


$D_\Gamma(\Omega) \subset W_r^2(\Omega) \bigcap \ker \gamma$, то


по теореме \ref{C0densV0} множество $W_r^2(\Omega) \bigcap \ker


\gamma$ плотно в $\overset{\circ}{V}$. Итак, осталось доказать


плотность множества $\gamma(W_r^2(\Omega))$ в


$\gamma(V)=W_q^1(\Gamma) \bigcap W_p^{1-1/p}(\Gamma)$.


Проведем следующие построения.


Поскольку в угловых точках множества $\Gamma$ гладкие куски $\Gamma$


стыкуются под ненулевым углом, то найдется конечное семейство ${\cal O}


=\{O_j\}_{j=1,N}$


попарно непересекающихся открытых в $R^2$ множеств, удовлетворяющее


следующим условиям:





1) $\Gamma \subset \bigcup \limits_{j=1}^{N} \overline{O_j}$;





2) $\Gamma \bigcap O_j \in C^2 \quad \forall O_j \in {\cal O}$.





\noindent{}Положим


$P=\bigcup\limits_{i \ne j} \overline{O_i}\cap \overline{O_j}\cap \Gamma$.


Очевидно, множество $P$ конечно и содержит в себе множество всех угловых


точек $\Gamma$ и все пересечения кривых $\Gamma_i$. Фиксируем


произвольную функцию $u \in V$.


Подберем функцию $u_0 \in C^\infty(\overline{\Omega})$ так, чтобы


$u_0(a) = u(a)$ $\forall a \in P$.


Тогда $v=u-u_0 \in V$ и функция $\gamma v=\gamma u - \gamma u_0$ обладает


тем свойством, что


$\gamma v(a) = 0$ $\forall a \in P$.


Пусть $\varepsilon > 0$ сколь угодно мало. Рассмотрим произвольное


множество $O_j \in {\cal O}$. Через $s$ обозначим дуговую координату


кривой $\overline{O_j} \cap \Gamma$, отсчитываемую от некоторой ее точки.


Функция $\gamma v$ как функция $s \in [s_1,s_2]$ принадлежит классу


$W_q^1\bigcap W_p^{1-1/p}(O_j\cap\Gamma)$. Рассмотрим два возможных случая.





1) Kусок $\overline{O_j} \cap \Gamma$ не содержит точек из $P$.


В этом случае кривая $\overline{O_j} \cap \Gamma$ 


замкнута и не имеет угловых точек. Поэтому найдется 


функция $\varphi _j \in C^\infty(\overline{O_j} \cap \Gamma)$


такая, что


$$


{\left\| {\gamma v-\varphi _j}


\right\|}_{W_q^1\cap W_p^{1-1/p}(O_j\cap\Gamma)}


\le \frac{\varepsilon }{N}.


$$


Так как $\overline{O_j} \cap \Gamma \in C^2$


и $\overline{O_j} \cap \Gamma \cap P = \emptyset$,


то найденную функцию $\varphi _j \in C^\infty(\overline{O_j} \cap \Gamma)$


можно продолжить с $\overline{O_j} \cap \Gamma$


на всю плоскость $R^2$ так, что


ее продолжение, которое обозначим через $v_j$, будет принадлежать


(по крайней мере) классу $C^2(R^2)$


и $\supp v_j \cap \overline{O_i}=\emptyset$ при $i \ne j$.





2) Kусок $\overline{O_j} \cap \Gamma$ содержит точки из $P$,


которые, очевидно, являются концами кривой\linebreak


$\overline{O_j} \cap \Gamma$.


В этом случае $\gamma v(s_1)=\gamma v(s_2)=0$. В силу известных


теорем о приближении финитными функциями функций


из соболевских классов, обращающихся в нуль на концах отрезка, найдется 


$\varphi _j \in C_0^\infty(O_j \cap \Gamma)$ (т.\,е. бесконечное число


раз дифференцируемая по $s$ функция, обращающаяся в нуль в окрестности


концов кривой $O_j \cap \Gamma$) такая, что


$$


{\left\| {\gamma v-\varphi _j}


\right\|}_{W_q^1\cap W_p^{1-1/p}(O_j\cap\Gamma)}


\le \frac{\varepsilon }{N}.


$$


Так как $\Gamma \bigcap O_j \in C^2$, 


то существует продолжение $w_j \in C^2(R^2)$


функции $\varphi _j \in C_0^\infty(O_j \cap \Gamma)$. Поскольку носитель


функции $\varphi_j$ целиком содержится в $O_j$, то найдется 


$\psi_j \in C_0^\infty(O_j)$, 


при этом $\psi_j=1$ в окрестности компакта $\supp \varphi _j$


(\cite{Her}, c.\,10).


Положим $v_j=w_j \psi_j \in C^2(R^2)$. 


По построению, $\supp v_j \cap \overline{O_i}


= \emptyset$ при $i \ne j$ (т.\,к. $O_j$ является окрестностью компакта


$\supp v_j$ и открытые множества $O_j$ и $O_i$ не перeсекаются)


и ${v_j|}_{O_j \cap \Gamma}=\varphi_j$.





Итак, для каждого $j=\overline{1,N}$ построены функции $v_j$ из $C^2(R^2)$,


удовлетворяющие условиям





1) $\supp v_j \cap \supp v_i = \emptyset$ при $j \ne i$;





2) $ {\left\| {\gamma v- \gamma v _j}


\right\|}_{W_q^1\cap W_p^{1-1/p}(O_j\cap\Gamma)} =


{\left\| {\gamma v-\varphi _j}


\right\|}_{W_q^1\cap W_p^{1-1/p}(O_j\cap\Gamma)}


\le \varepsilon /N$.





\noindent{}Положим


$v_\varepsilon = \sum \limits_{j=1}^{N} v_j |_\Omega$


и


$u_\varepsilon = v_\varepsilon + u_0$.


По построению $u_\varepsilon \in C^2(\overline{\Omega})


\subset W_r^2(\Omega)$ и


\begin{multline*}


{\left\| {\gamma u-\gamma u_\varepsilon}


\right\|}_{W_q^1\cap W_p^{1-1/p}(\Gamma)} =


{\left\| {\gamma v-\gamma v _\varepsilon}


\right\|}_{W_q^1\cap W_p^{1-1/p}(\Gamma)} \le


\sum \limits_{j=1}^{N} {\left\| {\gamma v-\gamma v_\varepsilon}


\right\|}_{W_q^1 \cap W_p^{1-1/p}(O_j\cap\Gamma)} = \\ =


\sum \limits_{j=1}^{N} {\left\| {\gamma v-\gamma v_j} \right\|}_{W_q^1


\cap W_p^{1-1/p}(O_j\cap\Gamma)} =


\sum \limits_{j=1}^{N} {\left\| {\gamma v-\varphi _j} \right\|}_{W_q^1


\cap W_p^{1-1/p}(O_j\cap\Gamma)} \le \varepsilon.


\end{multline*}


В силу произвольности $u \in V$ и $\varepsilon >0$ теорема доказана.


\medskip








Для $p_1,q_1\in [1,\infty)$ определим пространство функций $W$ переменных


$(t,x) \in Q=T \times \Omega$, где $T=(0,1)$, с конечной нормой


$$


{\left\| {u} \right\|}_W = {\bigg( {\int_{T} {\left\| {u(t)} \right\|}_


{W_p^1(\Omega)}^{^{p_1}} dt} \bigg)}^{_{1/p_1}} +


{\bigg( {\int_{T} {\left\| {\gamma u(t)} \right\|}_


{W_q^1(\Gamma)}^{^{q_1}} dt} \bigg)}^{_{1/q_1}}.


$$





\begin{thms} \label{CdensW}


Множество бесконечно дифференцируемых 


в $\overline{Q}=[0,1] \times \overline{\Omega}$


функций плотно в пространстве $W$.


\end{thms}





\begin{pf}


Как следует из определения, тип пространства $W$ тот же,


что и в предыдущем разделе, а именно $W=\gamma (L_{p_1}(T;W_p^1(\Omega)),


L_{q_1}(T;W_q^1(\Gamma)))$. Рассмотрим множество функций


$$


C^\infty[0,1] \otimes C^\infty(\overline{\Omega})=\left\{ {\varphi u:


\varphi \in C^\infty[0,1],


\ u \in C^\infty(\overline{\Omega})}\right\}\subset C^\infty (\overline{Q}).


$$


Как известно, множество $C^\infty[0,1]$ плотно в $L_{p_1}(0,1)


\cap L_{q_1}(0,1)=L_r(0,1)$, где $r=\max (p_1,q_1).$ По теореме


\ref{CdensV} \ $C^\infty(\overline{\Omega})$ плотно в


пространстве Соболева с усиленной метрикой


$V=\gamma(W_p^1(\Omega), W_q^1(\Gamma))$. Согласно следствию теоремы


\ref{dens} линейная оболочка $\mathrm{Span}\,(C^\infty[0,1]


\otimes C^\infty(\overline{\Omega})) \subset C^\infty


(\overline{Q})$ плотна в $W = \gamma (L_{p_1}(T;W_p^1(\Omega)),


L_{q_1}(T;W_q^1(\Gamma)))$.


\end{pf}








\begin{thebibliography}{30}





\bibitem{Timer}


Тимербаев М.Р. {\em Пространства с нормой графика и усиленные


пространства Соболева}. I //


Изв. вузов. Математика. -- 2003. -- \No\,5. -- C.\,55-65.





\bibitem{Danford}


Данфорд Н., Шварц Дж.Т. {\em Линейные операторы. Общая теория}.


-- М.: Ин. лит., 1962. -- 895~c.





\bibitem{IoffeLevi}


Иоффе А.Д., Леви В.Л. {\em Субдифференциалы выпуклых функций} //


Тр. Моск. матем. о-ва -- 1972. -- T.\,26. -- C.\,3--73.





\bibitem{Anton}


Антонцев С.Н., Мейрманов А.М. {\em Математические модели движения


поверхностных и подземных вод}. -- Новосибирск: Изд-во


Новосибирск. ун-та, 1979. -- 79~с.





\bibitem{GlazPav1}


Глазырина Л.Л., Павлова М.Ф. {\em Разностная схема решения задачи совместного


движения грунтовых и поверхностных вод} // Изв. вузов. Математика. --


1984. -- \No\,9. -- C.\,72--75.





\bibitem{GlazPav}


Глазырина Л.Л., Павлова М.Ф. {\em О разрешимости одного нелинейного


эволюционного неравенства теории совместного движения


поверхностных и подземных вод} // Изв. вузов. Математика. -- 1997. --


\No\,4. -- C.\,20--31.





\bibitem{Daut}


Даутов Р.З., Карчевский М.М., Федотов Е.М. {\em Об одном методе решения


трехмерных стационарных задач типа сосредоточенной емкости} // 


Tез. докл. Всесоюзн. конф. ``Актуальные проблемы вычислительной и


прикладной математики''. ВЦ СО АН СССР. -- Новосибирск, 1987. --


C.\,68.





\bibitem{Djak96}


D'akonov E.G. {\em Optimization in solving elliptic problems}. --


Boca Raton, 1996.





\bibitem{DjakDAN}


Дьяконов Е.Г. {\em Новый подход к краевым условиям Дирихле,


основанный на использовании усиленных пространств Соболева} //


Докл. РАН. -- 1997. -- T.\,352. -- \No\,5. -- C.\,590--594.





\bibitem{Djak}


Дьяконов Е.Г. {\em Усиленные пpостранства Соболева и некоторые новые типы


эллиптических краевых задач} // Дифференц. уравнения. -- 1997. -- T.\,33. --


\No\,4. -- C.\,532--539.





\bibitem{DjakIzVuz}


Дьяконов Е.Г. {\em Оценки $N$-поперечников в смысле Колмогорова для 


некоторых компактов в усиленных пространствах Соболева}


// Изв. вузов. Математика. --


1997. -- \No\,4. -- C.\,32--50.





\bibitem{Jak}


Яковлев Г.Н. {\em Граничные свойства функций класса $W_p^{(l)}$


на областях с угловыми точками} // ДАН СССР. -- 1961. -- T.\,140.


-- \No\,1. -- C.\,73--76.





\bibitem{Necas}


Necas J. {\em Les methodes directes en theorie des


equations elliptiques}. -- Masson, Paris~/~Academia, Pragua, 1967.


-- 346 p.





\bibitem{Her}


Хермандер Л. {\em Линейные дифференциальные операторы с частными 


производными}. -- М.: Мир, 1965. -- 380 с.





\end{thebibliography}





\vskip 2cm





\post{\it Казанский государственный}{\it Поступила}


\post{\it университет}{19.11.2001}





\label{end}


\end{document}


